%! AP Statistics — Unit 1, Lesson 8
\documentclass[10pt]{article}
\usepackage[margin=0.6in]{geometry}
\usepackage{xcolor}
\usepackage[most]{tcolorbox}
\usepackage{enumitem}
\usepackage{multicol}
\usepackage{amsmath,amssymb}
\usepackage{array}
\usepackage{tabularx}
\tcbuselibrary{breakable}

\definecolor{sky}{HTML}{EAF3FB}
\definecolor{navy}{HTML}{1F3A5F}
\definecolor{redbox}{HTML}{FCECEC}
\definecolor{greenbox}{HTML}{EDF8F0}
\definecolor{goldbox}{HTML}{FFF7E6}
\newtcolorbox{skillbox}[2][]{
    colback=#2, colframe=navy, boxrule=0.8pt, arc=2mm,
    left=2mm,right=2mm,top=1.5mm,bottom=1.5mm,
    fonttitle=\bfseries, title={#1}, halign title=center, breakable
}
\setlist[itemize]{leftmargin=1.2em,itemsep=0.35em,topsep=0.35em}
\setlist[enumerate]{leftmargin=1.4em,itemsep=0.35em,topsep=0.35em}
\newcolumntype{C}[1]{>{\centering\arraybackslash}p{#1}}
\newcolumntype{Y}{>{\centering\arraybackslash}X}

\begin{document}
\def\CourseName{AP Statistics}
\def\SchoolYear{2026--2027}
\def\MeetingLength{55 minutes}
\def\UnitNumberName{Unit 1: Exploring One-Variable Data \quad}
\def\LessonNumberName{Lesson 1.8: Graphical Representations of Summary Statistics}

\begin{center}
    {\Large\bfseries \CourseName: \SchoolYear\ (\MeetingLength) \\
    {\small\bfseries \UnitNumberName \LessonNumberName}}
\end{center}

\begin{tcolorbox}[colback=sky,colframe=navy,boxrule=0.9pt,arc=2mm,
    left=2mm,right=2mm,top=1.5mm,bottom=1.5mm]
    \textbf{Primary Objective:} Students will construct and interpret boxplots (box-and-whisker
    plots) from the five-number summary. Students will identify outliers on a modified boxplot and
    compare the information provided by a boxplot versus a histogram (Big Idea: UNC).
\end{tcolorbox}

\begin{skillbox}[Priority Ideas \& Skills]{goldbox}
\begin{minipage}[t]{0.48\linewidth}
    \textbf{AP Skill 2 --- Data Analysis}
    \begin{itemize}
        \item Construct a boxplot (including modified boxplot with outliers plotted separately) from a five-number summary.
        \item Read a boxplot to estimate median, IQR, range, and identify outliers.
        \item Describe a distribution from a boxplot using SOCS.
    \end{itemize}
\end{minipage}
\hfill
\begin{minipage}[t]{0.48\linewidth}
    \textbf{Key Understandings}
    \begin{itemize}
        \item A boxplot displays the five-number summary visually; each section contains approximately 25\% of the data.
        \item Boxplots do not show shape as clearly as histograms --- they hide gaps, clusters, and multimodality.
        \item Modified boxplots show outliers as individual points with whiskers extending to the last non-outlier.
    \end{itemize}
\end{minipage}
\end{skillbox}

\begin{skillbox}[{Vocabulary, Concepts, \& Theorems}]{greenbox}
\begin{minipage}[t]{0.48\linewidth}
    \begin{itemize}
        \item \textbf{Five-number summary}: $\min,\; Q_1,\; \text{Median},\; Q_3,\; \max$
        \item \textbf{Boxplot (box-and-whisker plot)}: box spans $Q_1$ to $Q_3$; line inside box at median; whiskers extend to $\min$ and $\max$ (or last non-outlier in modified version)
        \item \textbf{Modified boxplot}: whiskers extend to last non-outlier; outliers plotted as individual points ($\bullet$)
    \end{itemize}
\end{minipage}
\hfill
\begin{minipage}[t]{0.48\linewidth}
    \begin{itemize}
        \item \textbf{IQR box}: the central box; spans the middle 50\% of data
        \item \textbf{Whisker}: line from $Q_1$ to lower fence (or $\min$) and from $Q_3$ to upper fence (or $\max$)
        \item \textbf{Fence}: $Q_1 - 1.5\,\text{IQR}$ (lower) and $Q_3 + 1.5\,\text{IQR}$ (upper); values beyond fences are outliers
        \item \textbf{Percentile}: the value below which a given percentage of data falls
    \end{itemize}
\end{minipage}
\end{skillbox}

\begin{skillbox}[Activate Prior Knowledge \& Spiral Review]{sky}
\begin{minipage}[t]{0.48\linewidth}
    \textbf{Quick Review (3 min)}\\
    Students compute the five-number summary for a small dataset (5--7 values) from memory. Connect: today we draw those five numbers on a number line in a structured way called a boxplot.
\end{minipage}
\hfill
\begin{minipage}[t]{0.48\linewidth}
    \textbf{Spiral Connections}
    \begin{itemize}
        \item Lesson 1.7: five-number summary and IQR outlier rule needed to construct a modified boxplot.
        \item Lesson 1.9: side-by-side boxplots are the primary tool for comparing two distributions.
        \item AP exam: boxplots appear frequently in free-response questions requiring comparison.
    \end{itemize}
\end{minipage}
\end{skillbox}

\begin{skillbox}[Hook \& Lesson]{sky}
\begin{minipage}[t]{0.48\linewidth}
    \textbf{Hook (4--5 min)}\\
    Display a boxplot with no labels. Ask: ``What does each part mean? What can you determine just by looking?'' Students identify the box and whiskers intuitively before the formal introduction. Then: ``What \textit{can't} you tell from a boxplot?'' (Gaps, modes, exact values.) This motivates when to use each display.
\end{minipage}
\hfill
\begin{minipage}[t]{0.48\linewidth}
    \textbf{Lesson Flow (15 min)}
    \begin{enumerate}
        \item Model constructing a boxplot from a five-number summary step-by-step.
        \item Identify an outlier using the fence; construct a modified boxplot.
        \item Read information from a boxplot: median, IQR, presence of outliers, approximate spread.
        \item Place a boxplot and histogram of the same data side by side; compare what each shows.
    \end{enumerate}
\end{minipage}
\end{skillbox}

\begin{skillbox}[Explicit Instruction \& Active Monitoring]{sky}
\begin{minipage}[t]{0.48\linewidth}
    \textbf{Direct Instruction (10 min)}
    \begin{itemize}
        \item Draw each construction step on the board: number line $\to$ mark five numbers $\to$ draw box $\to$ draw median line $\to$ draw whiskers $\to$ plot outliers.
        \item Emphasize: in a modified boxplot, the whisker ends at the last \textit{non-outlier} value, not at the fence.
        \item Show how a right-skewed distribution produces a longer right whisker (or upper outliers).
    \end{itemize}
\end{minipage}
\hfill
\begin{minipage}[t]{0.48\linewidth}
    \textbf{Active Monitoring}
    \begin{itemize}
        \item Check: Is the median line inside the box (not at an edge)? Are outliers plotted as dots, not included in the whisker?
        \item Probe: ``Which box section is widest? What does that tell you about the data?''
        \item Common error: extending the whisker to the fence value rather than the last non-outlier data point.
    \end{itemize}
\end{minipage}
\end{skillbox}

\begin{skillbox}[Group Work \& Differentiation]{redbox}
\begin{minipage}[t]{0.48\linewidth}
    \textbf{Group Activity (8--10 min)}\\
    Groups receive a dataset of 15 values. Each group:
    \begin{enumerate}
        \item Computes the five-number summary and checks for outliers.
        \item Constructs a modified boxplot (ruler required).
        \item Writes a SOCS description from the boxplot alone.
        \item Discusses: ``What would a histogram of this data look like? Sketch a prediction.''
    \end{enumerate}
\end{minipage}
\hfill
\begin{minipage}[t]{0.48\linewidth}
    \textbf{Differentiation}
    \begin{itemize}
        \item \textit{Support}: Provide the five-number summary pre-computed; students focus on construction and interpretation.
        \item \textit{Extension}: Given only a boxplot image, students reverse-engineer the approximate five-number summary and write a paragraph comparing two boxplots.
        \item \textit{Technology}: Verify boxplot with graphing calculator (1-Var Stats + BoxPlot in STAT PLOT).
    \end{itemize}
\end{minipage}
\end{skillbox}

\begin{skillbox}[Individual Work \& Assessment]{redbox}
\begin{minipage}[t]{0.48\linewidth}
    \textbf{Exit Ticket (5 min)}\\
    Given five-number summary: $\min=10$, $Q_1=18$, $M=25$, $Q_3=34$, $\max=60$. IQR = 16. Students: (1) check if 60 is an outlier, (2) construct the modified boxplot, (3) estimate what percentage of data falls between 18 and 34.
\end{minipage}
\hfill
\begin{minipage}[t]{0.48\linewidth}
    \textbf{AP-Style Check}\\
    \textit{``Based on the boxplot, describe the distribution of delivery times. Include all relevant features.''}\\[4pt]
    Students must address SOCS with context. Key AP expectation: do not just list numbers --- interpret them (e.g., ``The median delivery time was 25 minutes, meaning half of all deliveries took less than 25 minutes.'')
\end{minipage}
\end{skillbox}

\begin{skillbox}[Reinforcement \& Extension]{goldbox}
\begin{minipage}[t]{0.48\linewidth}
    \textbf{Homework / Practice}
    \begin{itemize}
        \item Construct a modified boxplot from two provided datasets; describe each using SOCS.
        \item Multiple-choice: reading values from a boxplot (median, IQR, outliers, spread).
    \end{itemize}
\end{minipage}
\hfill
\begin{minipage}[t]{0.48\linewidth}
    \textbf{Extension / Preview}
    \begin{itemize}
        \item \textit{Extension}: Compare a boxplot and histogram of the same data; list 3 things visible in one that are not visible in the other.
        \item \textit{Preview}: Lesson 1.9 places two or more boxplots side by side. Think about what features you would compare (center, spread, outliers, skewness) when two groups are shown together.
    \end{itemize}
\end{minipage}
\end{skillbox}

\end{document}
